\documentclass[../main.tex]{subfiles}

\begin{document}
\section{Random Calculations}
\subsection{The case $\left( n = 3 \right)$}
This case is the just the classic Golden Ratio calculation. 
Explicitly, 
\[
  f_3\left( x \right) = x^2 - x - 1 = 0
\]
its roots are given by
\[
  \alpha = \frac{ 1 + \sqrt{5} }{ 2 }, \quad 
  \beta = \frac{ 1 - \sqrt{5} }{ 2 }
\]
The difference between them is $\sqrt{5}$ which indeed generates the unique prime ideal above $5$ in $\O_{\mathbb{Q}(\sqrt{5})}$ and thus the unique element of $G\left( f_3 \right)$ is ramified above $5$ as expected.
Explicitly, both solutions are mapped to $3 \in \F_5$ and hence the unique of $G\left( f_3 \right)$ acts trivially on the residue field.

In a more modern (but maybe in this case redundant) language, we could note that $\alpha$ generates $\O_{\Q_5(\sqrt{5})}$ over $\O_{\Q_5}$.
Indeed, following the procedure presented in Serre's proof of \cite[III \S 6 Proposition 12]{SerreLocalFields}, we can note that 
\[
  v_5\left( \alpha + 2 \right) = 
  v_5\left( \frac{ 5 + \sqrt{5} }{ 2 } \right) = 1
\]
and hence 
\[
  1, \alpha, \alpha + 2, \alpha^2 + 2\alpha
\]
generate $\O_{\Q_5(\sqrt{5})}$ over $\O_{\Q_5}$ as a module [maybe I need to show it explicitly at some point].
Hence the fact that the unique element of $G\left( f_3 \right)$ maps $\alpha$ to $\beta$ which is in the same $\left( \sqrt{5} \right)$ coset, implies that it belongs to the ramification group.

\subsection{The case $\left( n = 5 \right)$}
First off, a direct calculation shows that 
\[
  d_5 = \pm 563
\]
in particular, it ramifies only above $563$.
Moreover, in $\F_{563}$, the unique double root is 
\[
  \alpha := \frac{2 \cdot 4}{5} \equiv 
  8 * 338 \equiv 452 \mod 563
\]

\subsection{The Case $n = 10$}
Let $L$ be the splitting field of $f_{10}$ over $\Q_2$.
\begin{lemma}
  $L$ is ramified over $\left( 2 \right)$.
\end{lemma}
\begin{proof}

  It suffices to show that $\O_L$ contains an element of fractional valuation.
  Note that over $\F_2$,
  \[
    \overline{g_{10}}\left( x \right) = x^{10} + 1 = \left( x^5 + 1 \right)^2 = \left( x + 1 \right)^2 \left( x^4 + x^3 + x^2 + x + 1 \right)^2 
  \]
  Therefore, the roots of $g_{10}$ over $\Q_2$ are given by two lifts of $1$ and two lifts of each root of $x^4 + x^3 + x^2 + x + 1$ over $\F_{2^4}$.
  Choose $\overline{\beta}$ to be a root of $x^4 + x^3 + x^2 + x + 1$ in $\F_{2^4}$.
  Let $\beta \in \O_{L}$ be a lift of $\overline{\beta}$ (not necessarily a root of $g_{10}$).
  Consider the polynomial, 
  \[
    f_{10}\left( x + \beta \right)
  \]
  Its roots are given by $\alpha_i - \beta$ where $\alpha_i$ are the roots of $f_{10}$. 
  Therefore, if we fix a prime $\P \subseteq \O_L$ lying over $(2)$, by Lemma \ref{lem:structure-of-roots}, $f_{10}\left( x + \beta \right)$ has exactly two roots in $\P$.
\end{proof}

We will show that $G_{10}$ is primitive.
Let $\alpha_1, \ldots, \alpha_9$ be the roots of $f_{10}$.
Then in $\F_2$,
\begin{align*}
  & \overline{\left( x - 1 \right) \prod_{i=1}^{9} \left( x - \alpha_i \right)} = \\
  & \overline{g_{10}}\left( x \right) = \\
  & x^{10} - 2x^9 + 1 = \\
  & x^{10} + 1 = \\
  & \left( x^5 + 1 \right)^2 = \\
  & \left( x + 1 \right)^2 \left( x^4 + x^3 + x^2 + x + 1 \right)^2 
\end{align*}

\subsection{The case $n = 2p$ for $p$ an odd prime}

We will show that $G_n$ is primitive.
Assume towards contradiction that there is a partition of the $n - 1$ roots of $f_n$ into blocks respceted by $G_n$.
Since $G_n$ acts transitively on the roots, all blocks must be of the same size, say $d$, and denote the blocks by $B_1, \ldots, B_k$ where $k = \frac{ n - 1 }{ d }$.

Note now that over $\F_2$, $g_n$ factors as
\[
  g_n = x^n + 1 = \left( x^p + 1 \right)^2 = 
  \left( x + 1 \right)^2 \left( \sum_{i=0}^{p-1} x^i \right)^2
\]
and hence 
\[
  \overline{f_n} = \overline{\frac{ g_n }{ x - 1 }} = 
  \left( x + 1 \right) \left( \sum_{i=0}^{p-1} x^i \right)^2
\]

Fix a prime $Q$ of $\O_{L_n}$ lying over $(2)$.
Let $\alpha$ be the unique root of $f_n$ that projects to $1$ and assume WLOG that $\alpha \in B_1$.
Since it forms its own decomposition block, the decomposition group at $Q$ stabilizes $\alpha$.
In particular, it stabilizes $B_1$.
Therefore, since by lemma [will be added later], the inertia group at $Q$ acts transitively on each inertia block, the rest of $B_1$ is a union of $\frac{d - 1}{2}$ inertia blocks.
Therefore, since the frobenious automorphism at $Q$ acts on the set of inertia blocks as a cycle of length $\ord_p\left( 2 \right)$, it follows that $\ord_p\left( 2 \right) \mid \frac{d - 1}{2}$.

Now, fix $1 < i \leq k$.
Since $\#B_i = d$ is odd, $B_i$ cannot be a union of inertia blocks, and hence it must contain at least one root, $\alpha$ such that there is at least one automorphism $\sigma$ that maps $\alpha$ to a root in $B_j$ for some $j \neq i$.
However, since $G$ respects the block structure, $\sigma$ must map the entire block $B_i$ to $B_j$.
Therefore, for each root $\beta \in B_i$, the unique other root that shares its inertia block must also be in $B_j$.

Let $\phi \in D\left( \P \right)$ be the frobenious.
Then one of two options must hold:
\begin{enumerate}
  \item $\phi\left( B_i \right) = B_i$. But this implies that $\ord_p\left( 2 \right) \mid d$ which is impossible since then $\ord_p\left( 2 \right) \mid \gcd\left( d, \frac{d - 1}{2} \right) = 1$.
  \item There is $\alpha \in B_i$ such that $\phi\left( \alpha \right) \in B_j$ for some $j \neq i$.
\end{enumerate}


and denote by $I \leq G_n$ the inertia group 
corresponding to $Q$.
Then the map $\pi_Q$ reduces to a surjective $2$ to $1$ map from the roots of $g_n$ to the $\mu_{p - 1} \subseteq \overline{\F_{2}}$ and the action of $I$ stabilizes the fibers of this map.

In particular, the set of roots of $g_n$ includes exactly two lifts of $1$ which the action of $I$ either fixes or swaps.
However, $1 \in \Q$ must be one of these two lifts, and hence $I$ must fix both of them.
Denote the second lift of $1$ by $\alpha$ and denote by $B_1$ the block containing $\alpha$.
Then since $G_n$ respects the block structure, $I$ must stabilize $B_1$.

By lemma [will be added later], the action of $I$ on each inertia block is transitive.
Therefore,

\subsection{The case $n = 2^d$}


\end{document}